\section{Funkce více proměnných I}
\begin{oznaceni}
    \,
    \begin{itemize}
	\item $\mathbb R ^n = \underbrace{\mathbb R \times \mathbb R  \times \dots \times \mathbb R}_{k\ \text{krát}} $
	\item $x \in \mathbb R ^n$ znamená, že $x=[x_1,x_2,\dots,x_n]$
	\item $e^i = [0,0,\dots,0, \underset{i-\text{tý}}{1},0,\dots,0]$
	\item $o=[0,\dots,0]$
	\item $\nm{x}= \sqrt{\sum_{i=1}^\infty x_i^2}$
	\item někdy píšeme $\rho_e(x,y)=\nm{x-y}$
    \end{itemize}
\end{oznaceni}

\begin{definice}
Nechť $f$ je reálná funkce $n$ proměnných,
$a\in\mathbb{R}^n$
a $1\le i\le n$.
Pak \textbf{parciální derivaci funkce $f$ v bodě $a$
podle $i$-té proměnné} definujeme jako limitu

\[
\frac{\partial f}{\partial x_i}(a)
=
\lim_{t\to0}
\frac{
f(a+t e^i)-f(a)
}{t}.
\]
Symbolem
$
\frac{\partial f}{\partial x_i}
$
označujeme \textbf{parciální derivaci funkce $f$
podle $i$-té proměnné}, tj. funkci definovanou
předpisem
\[
\frac{\partial f}{\partial x_i}
:
x\mapsto
\frac{\partial f}{\partial x_i}(x).
\]
\end{definice}

\begin{poznamka}\,
\begin{enumerate}[(1)]
	\item Značení: místo symbolu $\frac{\partial f }{\partial x_i}(a)$ někdy používáme symbol $\partial_i f(a)$.
	\item Limita v definici
	$$
 \begin{cases}
		\text{neexistuje},\\
		\text{existuje} \ \begin{cases}
			\in \mathbb R \ \text{vlastní},\\
			\text{nevlastní} \ \begin{cases}
				\infty,\\
				-\infty.
\end{cases}
\end{cases}
 \end{cases}
	$$
	\item Pokud $\frac{\partial f }{\partial x_i}(a)$ existuje, pak
	$$
	    \exists \delta>0: \left \{ a+te^i, |t|<\delta \right \}\subset \mathcal D(f).
	$$
	\item Početní technika. Nechť $f$ je funkce $n$ proměnných, $a \in \mathbb R ^n,$ $1\leq i \leq n$.
	Parciální funkce je definována jako
	$$
	g(x)=f(a_1,\dots,a_{i-1},x,a_{i+1},\dots,a_n).
	$$
	Fukce $g$ je z $\mathbb R $ do $\mathbb R $.
	Platí
    $$
    g^\prime(a)=\frac{\partial f }{\partial x_i}(a),
    $$
    pokud má jedna strana smysl.
\end{enumerate}
\end{poznamka}

\begin{priklad}
	Vezměme funkci $f(x_1,x_2)=e^{x_1x_2^2}$. Spočteme parciální derivace
	\begin{align*}
		\frac{\partial f }{\partial x_1} (x_1,x_2) &= e^{x_1x_2^2}x_2^2,\\
		\frac{\partial f }{\partial x_2} (x_1,x_2) &= e^{x_1x_2^2}2x_1x_2.
\end{align*}
\end{priklad}

\begin{definice}
Nechť $f$ je reálná funkce $n$ proměnných,
$a\in\mathbb{R}^n$
a $L:\mathbb{R}^n\to\mathbb{R}$ je lineární zobrazení.
Řekneme, že $L$ je \textbf{totální diferenciál funkce
$f$ v bodě $a$}, jestliže platí
\[
\lim_{h\to0}
\frac{
f(a+h)-f(a)-L(h)
}{
\nm{h}
}
=
0.
\]
\end{definice}

\begin{poznamka}
	Pro $n=1$ máme
	$$
 \frac{f(a+h)-f(a)-Ah}{|h|} \to 0.
	$$
	Tento výraz můžeme ekvivalentně přepsat jako
	$$
 \frac{|f(a+h)-f(a)-Ah|}{|h|} \to 0.
	$$
	Jinými slovy,
	$$
	    \left |\frac{f(a+h)-f(a)-Ah}{h}\right | \to 0.
	$$
	Konečně
	$$
	    \left |\frac{f(a+h)-f(a)}{h} - A\right | \to 0.
	$$
	Proto $A=f^\prime(a)$.
	\par Pro $n=2$ máme
	$$
	    (x,y)\mapsto f(x,y).
	$$
	Lineární zobrazení bude
	$$
	    (x,y)\mapsto f(a_1,a_2) + L(x-a_1,y-a_2).
	$$
	\par Pro $n \in \mathbb N $ máme zobrazení $x \mapsto f(x)$ a lineární zobrazení je
	$$
	    x \mapsto f(a)+L(x-a).
	$$
	Pak
	$$
 \frac{f(x)-f(a)-L(x-a)}{\nm{x-a}}
	$$
	nám po jednoduché substituci dává definiční výraz.
\end{poznamka}

\begin{poznamka}
	Podmínka z definice ohledně aproximace je ekvivalentní podmínce
	$$
 \lim_{x\to a } \frac{f(x)-f(a)-L(x-a)}{\nm{x-a}}=0.
	$$
\end{poznamka}

\begin{veta}[Vztah totálního diferenciálu a parciální derivace]\label{veta10.1}
Nechť $L$ je totální diferenciál funkce
$f$ v bodě $a\in\mathbb{R}^n$.
Potom existují parciální derivace
\[
\frac{\partial f}{\partial x_1}(a),
\dots,
\frac{\partial f}{\partial x_n}(a)
\]
a pro každé $h\in\mathbb{R}^n$ platí
\[
L(h_1,\dots,h_n)
=
\frac{\partial f}{\partial x_1}(a)h_1
+\cdots+
\frac{\partial f}{\partial x_n}(a)h_n.
\]
\end{veta}

\begin{proof}
	Zobrazení $L$ pišme ve tvaru
	$$
	    L(h_1,h_2,\dots,h_n)=A_1h_1+ A_2h_2 + \dots + A_nh_n,
	$$
	kde $A_1,\dots,A_n \in \mathbb R $. Zvolme $i \in \left \{ 1,\dots,n \right \} $.
	Zobrazení $\varphi(t)=te^i$ je spojité. Pak je
	$$
	    \nm{\varphi(t)-\varphi(0)}=\nm{te^i-0e^i} = \nm{ [ 0,\dots, 0, \underset{i-\text{tý}}{t},0\dots,0 ] }=|t|.
	$$
	Potom podle věty \ref{veta9.10} platí
	$$
 \lim_{t\to 0} \frac{f \left( a+\varphi(t) - f(a) -L(\varphi(t)) \right ) }{\nm{\varphi(t)}} = 0,
	$$
	podmínka (P) je splněna. Upravíme
	$$
 \lim_{t\to 0} \frac{f(a+te^i)-f(a)-A_it}{|t|}=0.
	$$
	Odsud dostáváme, že
	$$
 \lim_{t \to 0} \left | \frac{f(a+te^i)-f(a)}{t}-A_i \right | = 0,
	$$
	a tedy
	$$
	\lim_{t \to 0} \frac{f(a+te^i) - f(a)}{t} = A_i.
	$$
	Proto $\frac{\partial f }{\partial x_i}(a)=A_i$.
\end{proof}

\begin{poznamka}
	Z věty \ref{veta10.1} plyne, že totální diferenciál je určen jednoznačně.
\end{poznamka}

\begin{oznaceni}
	Totální diferenciál značíme $f^\prime(a)$. Formálně je $f^\prime(a)\in \left( \mathbb R ^n \right )^* $.
	\par Pozor na kolizi značení! Pro $n=1$ je
	\begin{align*}
		f^\prime(a)\in \mathbb R && f^\prime(a) \in (\mathbb R)^* : \mathbb R \to \mathbb R
	\end{align*}
	třeba rozlišovat mezi derivací a totálním diferenciálem.
\end{oznaceni}

\begin{veta}\label{veta10.2}
Má-li funkce $f$ v bodě
$a\in\mathbb{R}^n$
totální diferenciál, je
$f$ v bodě $a$ spojitá.
\end{veta}

\begin{proof}
    Stačí ukázat, že následující limita je rovna $f(a)$:
    $$
    \lim_{x \to a } f(x) = \lim_{x \to a }  \left (\frac{f(x)-f(a)-f^\prime(a)(x-a) }{\nm{x-a}} \cdot \nm{x-a} + f(a) + f^\prime(a)(x-a)\right ),
    $$
    kde $f^\prime(a)$ je zobrazení z $\mathbb R ^n$ do $\mathbb R $, $h \mapsto f^\prime(a)(h)$.
    Protože je $x \mapsto \nm{x-a}$ vlastně
    $$x \mapsto \sqrt{\sum_{i=1}^n (x_i-a_i)^2},$$ je spojité.
    Konečně třetí člen je vlastně $$x \mapsto x-a = [x_1-a_1, \dots, x_n-a_n],$$ a tedy je též spojité.
    Po substituci je $$y \to f^\prime(a)(y) = \sum_{i=1}^n \frac{\partial f }{\partial x_i}(a)\cdot y_i.$$
    Máme tedy
    $$
    \lim_{x \to a } f(x) = 0\cdot 0 + f(a) + f^\prime(a) (\vec o) = f(a),
    $$
    což jsme chtěli dokázat.
\end{proof}

\begin{poznamka}
	\,
	\begin{enumerate}[(1)]
		\item Existuje-li $f^\prime(a)$, již jsme dokázali, že $\frac{\partial f }{\partial x_i}(a)\in \mathbb R $
		a také že $f$ je v $a$ spojité. Opačná implikace neplatí.
		\item Uvažujme
		$$
		    f(x,y) = \begin{cases}
							0 & x=0 \lor y=0, \\
							1 & \text{jinak}.
						\end{cases}
		$$
		Potom
		\begin{align*}
		\frac{\partial f }{\partial x }(0,0) = 0, & & \frac{\partial f }{\partial y}(0,0) = 0,
\end{align*}
ale $f$ není spojitá v $(0,0)$.
	\end{enumerate}
\end{poznamka}

\begin{lemma}\label{lemma10.3}
Nechť $f$ je reálná funkce $n$ proměnných,
$
I=
(\alpha_1,\beta_1)
\times\cdots\times
(\alpha_n,\beta_n)
\subset\mathbb{R}^n,
$
$a,b\in I$.
Nechť v každém bodě $I$ existují vlastní
parciální derivace $f$
podle všech proměnných.
Potom existují body
$
\xi_1,\dots,\xi_n\in I
$
takové, že
\[
f(b)-f(a)
=
\sum_{i=1}^{n}
\frac{\partial f}{\partial x_i}
(\xi_i)
(b_i-a_i).
\]
\end{lemma}

\begin{proof}
	Označme
	\begin{align*}
	p^0 &= (a_1,a_2,\dots,a_n) = a,\\
	p^1 &= (b_1,a_2,\dots, a_n),\\
	p^2 &= (b_1,b_2, \dots, a_n),\\
	&\dots \\
	p^n &= (b_1,b_2,\dots,b_n) = b.
\end{align*}
Potom platí
$$
    f(b) - f(a) = \sum_{k=1}^n \left( f(p^k)-f(p^{k-1}) \right ).
$$
Položme pro $k\in \left \{ 1,\dots,n \right \} $
$$
    g^k(x)=f(b_1, b_2,\dots, b_{k-1}, x, a_{k+1}, \dots, a_n).
$$
Potom má $g^k$ v každém bodě $(\alpha_k, \beta_k)$ vlastní derivaci, takže existuje $z_k \in [a_k, b_k]$
splňující (Lagrangeova věta \ref{veta5.7})
$$
    \underbrace{g^k(b_k)-g^k(a_k)}_{f(p^k)-f(p^{k-1})}=\underbrace{(g^k)^\prime(z_k)(b_k-a_k)}_{\frac{\partial f }{\partial x_k}(\xi^k)\cdot (b_k-a_k)}.
$$
Položíme $\xi^k = [b_1,b_2,\dots,b_{k-1},z_k, a_{k+1}, \dots,a_n]$.
\end{proof}

\begin{veta}\label{veta10.4}
Nechť $f$ je reálná funkce
$n$ proměnných,
$a\in\mathbb{R}^n$
a
\[
\frac{\partial f}{\partial x_1},
\dots,
\frac{\partial f}{\partial x_n}
\]
jsou spojité funkce v bodě $a$.
Potom má $f$ v bodě $a$
totální diferenciál.
\end{veta}

\begin{proof}
	Máme funkci $f:\mathbb R ^n \to \mathbb R $ a spolu s ní $n$-tici funkcí
	\begin{align*}
	\frac{\partial f }{\partial x_1} &: B(\alpha_1,\beta_1) \to \mathbb R , \\
	\frac{\partial f }{\partial x_2} &: B(\alpha_2,\beta_2) \to \mathbb R , \\
	& \dots\\
\frac{\partial f }{\partial x_n} &: B(\alpha_n,\beta_n) \to \mathbb R .
	\end{align*}
	Dále máme $L: \mathbb R ^n \to \mathbb R $ takové, že
	$$
	    L: k \mapsto \sum_{i=1}^n \frac{\partial f }{\partial x_i}(a)k_i.
	$$
	Ukážeme, že $L = f^\prime(a)$, tj. pro $L$ platí
	$$
 \lim_{x \to a} \frac{f(x)-f(a)-L(x-a)}{\nm{x-a}} = 0.
	$$
	Zvolme $\varepsilon>0$. K němu nalezneme $\tilde \delta >0$ takové, že
	$$
	    \forall i \in \left \{ 1, \dots, n \right \} \, \forall x \in B(a,\tilde \delta): \left | \frac{\partial f }{\partial x_i }(x) - \frac{\partial f }{\partial x_i} (a) \right |  <\varepsilon.
	$$
	Položme $\delta =\frac{\tilde \delta }{\sqrt n}$. Pak
	$$
	    I = (a_1-\delta, a_1+\delta)\times \dots \times (a_n-\delta,a_n+ \delta) \subset (a,\tilde,\delta).
	$$
	Zvolme $x \in B(a,\delta) \subseteq I \subseteq B(a,\tilde \delta)$. Podle lemmatu
	\ref{lemma10.3} nalezneme body $\xi^i (x)\in I$, $i =1,\dots,n $ takové, že
	$$
	    f(x)-f(a) = \sum_{i=1}^n \frac{\partial f }{\partial x_i} (\xi^i (x))(x_i-a_i).
	$$
	Pak máme
	\begin{align*}
	|f(x)-f(a)-L(x-a)|&=\left | \sum_{i=1}^n \frac{\partial f }{\partial x_i} (\xi^i(x))(x_i-a_i) - \sum_{i=1}^n \frac{\partial f }{\partial x_i}(a)(x_i-a_i) \right | \\
	&= \left | \sum_{i=1}^n \left( \frac{\partial f }{\partial x_i}(\xi^i(x)) - \frac{\partial f }{\partial x_i}(a) \right ) \cdot (x_i-a_i) \right |\\
	& \leq \sum_{i=1}^n \underbrace{\left | \frac{\partial f }{\partial x_i} (\xi^i(x))-\frac{\partial f }{\partial x_i}(a) \right |}_{< \, \varepsilon}  \cdot \underbrace{|x_i-a_i|}_{\leq\,\nm{x-a}}\\
	& \leq n\cdot  \varepsilon\cdot  \nm{x-a}.
\end{align*}
	Odsud plyne, že
	$$
 \lim_{x \to a } \frac{f(x)-f(a)-L(x-a)}{\nm{x-a}} = 0.
	$$
\end{proof}

\begin{veta}[Řetízkové pravidlo]\label{veta10.5}
Nechť funkce
$
f_1,\dots,f_k
$
z $\mathbb{R}^n$ do $\mathbb{R}$
mají v bodě
$a\in\mathbb{R}^n$
totální diferenciál
a funkce $g$ z $\mathbb R ^k$ do $\mathbb R $
má v bodě
$
b=
(f_1(a),\dots,f_k(a))
$
totální diferenciál.
Definujme funkci
$h:\mathbb{R}^n\to\mathbb{R}$ předpisem
\[
h(x)
=
g(
f_1(x),
\dots,
f_k(x)
).
\]
Potom má $h$ v bodě $a$
totální diferenciál a pro
$
i\in\{1,\dots,n\}
$
platí
\[
\frac{\partial h}{\partial x_i}(a)
=
\sum_{j=1}^{k}
\frac{\partial g}{\partial y_j}(b)
\frac{\partial f_j}{\partial x_i}(a).
\]
\end{veta}
